\section{Methodology}
\label{sec:methodology}

We organise implementation into five phases (Table~\ref{tab:phases}), each with validation gates G1--G5.

\begin{table}[t]
  \centering
  \caption{Five-phase methodology summary with validation gates.}
  \label{tab:phases}
  \scriptsize
  \begin{tabular}{clp{3.8cm}p{3.2cm}}
    \toprule
    Ph. & Name & Primary outputs & Gate \\
    \midrule
    1 & Data foundation (O1) & STAC catalogue; \texttt{datasource\_id} & G1: ingest OK \\
    2 & Semantic model (O2) & TBox; alignment; seed OBDA & G2: ontology valid \\
    3 & LLM capability (O3--O4) & Dataset; fine-tuned model & G3: EX/EA thresholds \\
    4 & VKG automation (O5) & Draft \texttt{.obda} patches & G4: SQL+SPARQL pass \\
    5 & Integration (O6) & Provenance matrix; confidence & G5: case metrics \\
    \bottomrule
  \end{tabular}
\end{table}


\subsection{Phase 1: Data foundation (O1)}

Datasources (OSM extracts, census tiles, DTM/DSM rasters, grid assets, user uploads) are classified as spatial or non-spatial, registered via \texttt{POST /api/v1/datasources}, and ingested with manifests.
Each feature row carries \texttt{datasource\_id} for provenance back-links.
We measure ingest success rate, metadata completeness, provenance link rate, and ingest overhead versus direct load (H1).

\subsection{Phase 2: Semantic model (O2)}

The ontology stack has three layers:
\textbf{L1 Meta} (RDF, RDFS, OWL, SKOS, schema.org, Dublin Core);
\textbf{L2 Upper} (BFO, SUMO, DOLCE---planned);
\textbf{L3 Domain} (CityGML, BOT, SOSA/SSN, GeoSPARQL, OWL-Time, \texttt{cim:}).
File \texttt{alignment.ttl} asserts equivalences (e.g.\ \texttt{citygml:Building} $\equiv$ \texttt{bot:Building}) and subsumptions linking zones and observations.
Seed mappings in \texttt{cim\_citydb.obda} cover buildings, thermal zones, energy observations, geometries, and time series.

\subsection{Phase 3: LLM capability (O3, O4)}

\textbf{Dataset generation} (CIM Assist / \texttt{ai4db}).\\
Stage~1 (\texttt{stage1\_cim\_v2.py}) produced 10,800 rule-based question--SQL pairs across 52 templates (98--100\% NoErr).
Stage~2 augmented to 50,000 samples via GPT-4o-mini at \$50 total cost (97.34\% NoErr).
After multi-dimensional stratification (task type, domain, complexity, question tone), the training set contains \textbf{34,993} pairs (90\% positive spatial SQL; 10\% negative samples).
The held-out benchmark has \textbf{201} stratified samples with 100\% executable gold SQL~\cite{taherdoust2025cimassist}.

\textbf{Fine-tuning} (thesis; reused in Paper~2).\\
We fine-tuned \texttt{defog/sqlcoder-7b-2} with QLoRA~\cite{dettmers2023qlora,hu2022lora} on an NVIDIA RTX~6000 (IPAZIA cluster, Politecnico di Torino).
SQLCoder-7B completed training in 51\,hours (2,187 steps; evaluation loss 0.036) using 18\,GB VRAM.
Qwen~2.5-14B was also trained but rejected due to poor benchmark EX despite larger size~\cite{taherdoust2025cimassist}.
Table~\ref{tab:finetuning} summarises the production configuration.

\begin{table}[t]
  \centering
  \caption{QLoRA fine-tuning configuration (CIM Assist thesis; reused in O4/O5).}
  \label{tab:finetuning}
  \small
  \begin{tabular}{ll}
    \toprule
    Parameter & Value \\
    \midrule
    Base model & SQLCoder-7B (\texttt{defog/sqlcoder-7b-2}) \\
    Method & QLoRA 4-bit; rank 16; $\alpha$ 32 \\
    Training samples & 34,993 (90\% positive; 10\% negative) \\
    Effective batch size & 16 (batch 4 $\times$ grad.\ accum.\ 4) \\
    Learning rate & $2.0 \times 10^{-4}$ (cosine; 5\% warmup) \\
    Hardware & NVIDIA RTX 6000 24\,GB (IPAZIA, PoliTo) \\
    Training time & 51\,h; 2,187 steps; eval.\ loss 0.036 \\
    GPU memory & 18\,GB (75\% VRAM) \\
    \bottomrule
  \end{tabular}
\end{table}


\textbf{Evaluation.}
\texttt{evaluator\_v2.py} reports EM, EX, Deep EM, SC, and EA with taxonomy-stratified breakdowns (H4).
The fine-tuned SQLCoder-7B (V3) reaches \textbf{84.58\% EX} and \textbf{84.58\% EA} (89\% first-shot success; minimal gain from self-correction), versus 30.85\% EX for GPT-4o~\cite{taherdoust2025cimassist}.

\subsection{Phase 4: VKG automation (O5)}

\textbf{Pivot from CIM Assist.}
In the thesis, the fine-tuned model powered \texttt{agent\_cim\_assist.py} for interactive NL$\rightarrow$SQL over CIM Wizard.
In Paper~2, the \emph{same} model generates \emph{source SQL} for Ontop mappings; RDF \emph{target} templates are assembled from ontology rules (not free-form LLM output).

Inspired by LLM4VKG~\cite{xiao2025llm4vkg}, a multi-agent pipeline performs:
(1)~schema graph extraction;
(2)~ontology alignment with mapping-pattern recognition (SE, SR, SRm, SH);
(3)~source SQL generation via fine-tuned LLM;
(4)~target RDF assembly from template rules;
(5)~validation (SQL execution, Ontop load, SPARQL smoke tests);
(6)~human-approved merge into \texttt{.obda} (H5).

\subsection{Phase 5: Integration and confidence (O6)}

The \texttt{PipelineExecutor} supports \texttt{mode=all\_methods}.
For each $(\mathit{building}, \mathit{feature}, \mathit{method})$ we persist provenance:
method name, priority, status, value, input features, \texttt{datasource\_ids}, timestamp, and confidence level $c \in \{1,\ldots,5\}$.

Confidence combines weighted factors:
\textbf{source tier} (measured $>$ observed $>$ inferred $>$ default),
\textbf{method agreement} across the multi-method matrix,
\textbf{input completeness},
\textbf{datasource freshness},
and \textbf{execution status}.
Table~\ref{tab:confidence} defines illustrative ordinal levels.

\begin{table}[t]
  \centering
  \caption{Illustrative ordinal confidence levels (O6).}
  \label{tab:confidence}
  \small
  \begin{tabular}{clp{4.5cm}}
    \toprule
    Level & Label & Rule (illustrative) \\
    \midrule
    5 & Verified & Measured datasource + successful method \\
    4 & High & Tier-1 inferred; $\geq$2 methods agree within $\tau$ \\
    3 & Medium & Single tier-2 method; no conflict \\
    2 & Low & Default method or method conflict \\
    1 & Uncertain & Missing inputs; large cross-method spread \\
    \bottomrule
  \end{tabular}
\end{table}


Provenance and confidence are exported to the VKG via \texttt{prov:wasGeneratedBy} and \texttt{cim:confidenceLevel} properties, enabling SPARQL lineage queries (H6).

\begin{figure}[t]
  \centering
  \fbox{\parbox{0.92\linewidth}{\centering\vspace{1.5em}
    \textbf{Figure placeholder:} Five-phase methodology pipeline (Phases 1--5) with validation gates G1--G5.
  \vspace{1.5em}}}
  \caption{High-level methodology schema with validation gates.}
  \label{fig:methodology}
\end{figure}
